\documentclass[11pt]{article}
\usepackage[margin=0.9in,letterpaper]{geometry}
\usepackage{enumitem}
\usepackage{amsmath}
\usepackage{parskip}
\usepackage[colorlinks=true,linkcolor=black]{hyperref}
\setlist[enumerate]{itemsep=7pt, topsep=5pt}
\pagestyle{plain}
\begin{document}

\begin{center}
{\Large\textbf{Stat 220: Building and Trusting a Model}}\\[3pt]
{\large Midterm practice questions, with answers}\\[6pt]
\end{center}

Practice on many predictors, overfitting, honest evaluation, and how much flexibility your data
can actually pay for. No computer needed.

\vspace{8pt}
\hrule
\vspace{10pt}

\begin{enumerate}[itemsep=14pt]
\item A teammate adds 40 variables and training $R^2$ rises from 0.62 to 0.95. How impressed should you be?

\begin{enumerate}[label=\Alph*., itemsep=1pt, topsep=3pt]
  \item Very. That is a large improvement in fit.
  \item Not at all by itself. In-sample $R^2$ can only go up when you add variables, even useless ones, so the jump is expected rather than evidence.
  \item Somewhat. It depends on whether the 40 variables were significant.
  \item Not impressed, because 0.95 is suspiciously high and the data must be wrong.
\end{enumerate}

\vspace{2pt}
\textbf{Answer: B.} Adding a column can never reduce the fit on the data you fitted, so training $R^2$ is a ratchet. The number that would mean something is error on rows the model never saw. Ask for that, and ask how the 40 were chosen.
\vspace{4pt}

\item Why does the error curve have a minimum when neither of its two parts does?

\begin{enumerate}[label=\Alph*., itemsep=1pt, topsep=3pt]
  \item Because cross-validation introduces randomness.
  \item Because bias falls with complexity while variance rises, and only their sum turns around.
  \item Because the training error eventually starts rising.
  \item Because of the central limit theorem.
\end{enumerate}

\vspace{2pt}
\textbf{Answer: B.} Bias measures being wrong in the same direction every time, and it decreases the whole way as the model gains flexibility. Variance measures being different every refit, and it increases the whole way. Neither curve has a bottom. Their sum does, and finding it is what cross-validation is for.
\vspace{4pt}

\item What moves the best model complexity toward more flexibility?

\begin{enumerate}[label=\Alph*., itemsep=1pt, topsep=3pt]
  \item More noise in the outcome.
  \item More candidate predictors to search through.
  \item More rows of data.
  \item A stronger need to explain a coefficient.
\end{enumerate}

\vspace{2pt}
\textbf{Answer: C.} Data buys down variance, which is the cost of flexibility, so more rows make a complicated model affordable. The other three all push the other way. More noise means more of what you would fit is noise, searching more candidates is itself a form of fitting, and needing to defend a number rules out models whose parameters mean nothing.
\vspace{4pt}

\item Cross-validated errors for a depth-3 tree and a depth-7 tree are within one standard error of each other. What does the one-standard-error rule say?

\begin{enumerate}[label=\Alph*., itemsep=1pt, topsep=3pt]
  \item Take the depth-7 tree, since its mean was lower.
  \item Take the depth-3 tree, since the two are not distinguishable on this data.
  \item Try depths 4, 5 and 6 to break the tie.
  \item Collect more data.
\end{enumerate}

\vspace{2pt}
\textbf{Answer: B.} A cross-validated score is an estimate and carries noise of its own. Inside that noise, choosing the lower mean is choosing at random. The simpler model is more stable over time, easier to explain, and less likely to break under a shift in the data, and none of those are consolation prizes.
\vspace{4pt}

\item A model predicting hospital readmission achieves 0.99 AUC in testing. What is your first thought?

\begin{enumerate}[label=\Alph*., itemsep=1pt, topsep=3pt]
  \item Excellent, ship it.
  \item Suspicion. Check whether a variable in the model encodes the answer.
  \item The test set was too small.
  \item The model is underfitting.
\end{enumerate}

\vspace{2pt}
\textbf{Answer: B.} Near-perfect performance on a genuinely hard problem usually means leakage: something in the feature set is a consequence of the outcome rather than a predictor of it. A discharge code recorded at readmission, or a date field that only exists for readmitted patients. Real prediction problems are hard, and 0.99 is a claim that this one was not.
\vspace{4pt}

\item What is wrong with running cross-validation after scaling the entire dataset?

\begin{enumerate}[label=\Alph*., itemsep=1pt, topsep=3pt]
  \item Nothing, scaling is a standard preprocessing step.
  \item The scaling used the held-out rows, so information leaked from the test folds into training.
  \item Scaling makes coefficients uninterpretable.
  \item Cross-validation cannot be used with scaled data.
\end{enumerate}

\vspace{2pt}
\textbf{Answer: B.} The mean and standard deviation used for scaling were computed with the held-out rows included, so every fold's training set has already peeked at its test set. The effect is usually small but it is in the optimistic direction, and it is free to avoid by scaling inside the pipeline so it is refit on each fold's training portion.
\vspace{4pt}

\item Two predictors correlate at 0.95. What happens to their coefficients?

\begin{enumerate}[label=\Alph*., itemsep=1pt, topsep=3pt]
  \item Both become significant, since they share signal.
  \item The coefficients become unstable and each can look insignificant, even though together the pair clearly matters.
  \item The model refuses to fit.
  \item One is automatically dropped.
\end{enumerate}

\vspace{2pt}
\textbf{Answer: B.} The model cannot tell which twin deserves the credit, so it splits it in a way that swings wildly between refits, and the standard errors inflate. The tell is that each looks unimportant on its own while dropping both hurts badly. Insignificant here means indistinguishable from its twin, not unimportant.
\vspace{4pt}

\item You search 200 predictors against an outcome that is pure noise. Roughly what does the best one look like?

\begin{enumerate}[label=\Alph*., itemsep=1pt, topsep=3pt]
  \item Nothing, since the outcome is noise.
  \item Weakly correlated but not significant.
  \item Correlated around 0.4 and significant at $p<0.01$ if you report it without mentioning the search.
  \item Perfectly correlated, since one of 200 will match by chance.
\end{enumerate}

\vspace{2pt}
\textbf{Answer: C.} The maximum of 200 noisy correlations is a large number by construction. Reporting it as though it were the only thing you tried is the mechanism behind a great many findings that fail to replicate. If the variables were chosen by looking at the data, the printed $p$-values assume something that is not true.
\vspace{4pt}

\item What does stepwise selection by $p$-value do to the coefficients that survive?

\begin{enumerate}[label=\Alph*., itemsep=1pt, topsep=3pt]
  \item Nothing, they are unbiased.
  \item It shrinks them toward zero.
  \item It biases them away from zero, because surviving required looking large.
  \item It makes them all significant by construction.
\end{enumerate}

\vspace{2pt}
\textbf{Answer: C.} Survival was conditional on clearing a threshold, so the survivors are enriched for estimates that happened to land high. The reported intervals and $p$-values are computed as if the model had been fixed in advance, which it was not. Report out-of-sample performance instead, or validate on fresh data.
\vspace{4pt}

\item When is a simpler model the right choice even though it predicts worse?

\begin{enumerate}[label=\Alph*., itemsep=1pt, topsep=3pt]
  \item Never, prediction accuracy is the goal.
  \item When someone will have to defend an individual output, or maintain it for years.
  \item Only when the difference is under 1\%.
  \item When the dataset is small.
\end{enumerate}

\vspace{2pt}
\textbf{Answer: B.} Accuracy is one requirement among several. A model whose outputs get challenged by a regulator, a customer, or a colleague has to be explainable, and one that runs nightly for two years has to be stable. Choosing simplicity for a stated reason reads as judgment. Defaulting to it because you never tried anything else reads as inexperience, and the difference is whether you can say what you ruled out.
\vspace{4pt}

\end{enumerate}

\end{document}
